% Part: first-order-logic
% Chapter: lindstrom
% Section: lindstrom-proof

\documentclass[../../../include/open-logic-section]{subfiles}

\begin{document}

\olfileid{mod}{lin}{prf}

\olsection{د ليندستروم قضيه}

\begin{lem}
\ollabel{lem:lindstrom}
فرض کړئ~$!E \in L(\Lang{L})$ او~$\Lang{L}$ متناهي ده؛ دا هم فرض
کړئ چې داسې~$n \in \Nat$ شته چې د هر دوو !!{structure}s
$\Struct{M}$ او~$\Struct{N}$ دپاره، که~$\Struct{M} \equiv_n
\Struct{N}$ او~$\Struct{M} \models_L !E$ وي، نو
$\Struct{N} \models_L !E$ هم وي۔ بيا~$!E$ د لومړۍ درجې له يوې
!!{sentence} سره هم‌ارزه ده؛ يعنې د لومړۍ درجې داسې~$!D$ شته چې
$\Mod(L){!E} = \Mod(L){!D}$ وي۔
\end{lem}

\begin{proof}
$n$ داسې ونيسئ چې هر دوه~$n$-هم‌ارز !!{structure}s~$\Struct{M}$ او
$\Struct{N}$،~$!E$ ته د ورکړي ارزښت په اړه سره سلا وي۔
\olref[bas][pis]{prop:qr-finite} راياد کړئ: په يوه متناهي
!!{language} کښې، له منطقي هم‌ارزۍ پرته، د لومړۍ درجې يوازې متناهي
ډېرې داسې !!{sentence}s شته چې د کميت ټاکونکو رتبه ئې تر~$n$ زياته
نۀ وي۔ اوس د هر ټاکلي !!{structure}~$\Struct{M}$ دپاره
$!D_{\Struct{M}}$ د هغو د لومړۍ درجې !!{sentence}s~$!E$ د هر منطقي
هم‌ارزۍ له ټولګي څخه د يوه استازي تړون ونيسئ چې په~$\Struct{M}$ کښې
رښتونې وي او~$\QuantRank{!E} \le n$ ولري؛ دا تړون متناهي دے، او
$\Struct{N} \models !D_{\Struct{M}}$ هله او يوازې هله چې
$\Struct{N} \equiv_n \Struct{M}$ وي۔ بيا
$!D = \textstyle\bigvee
\Setabs{!D_{\Struct{M}}}{\Struct{M} \models_L !E}$ کېږدئ؛ له هر
منطقي هم‌ارزۍ ټولګي څخه د يوه استازي په ساتلو دا بېلتون هم متناهي
دے۔
\textbf{د ژباړې د نسخې سمون (OLMOD-028):} سرچينې د ټولو رښتونو
ټيټې مرتبې جملو تړون «متناهي» بللے و، که څه هم په هر منطقي هم‌ارزۍ
ټولګي کښې بې‌شمېره نحوي بڼې راتللے شي؛ دلته په څرګنده له هر ټولګي
څخه يو استازی اخيستل کېږي، او په وروستي بېلتون کښې هم تکراري هم‌ارزې
برخې لرې کېږي۔

پايله~$\Mod(L){!E} = \Mod(L){!D}$ تر لاسه کېږي۔ په حقيقت کښې، که
$\Struct{N} \models_L !D$ وي، نو د کوم
$\Struct{M} \models_L !E$ دپاره~$\Struct{N} \models
!D_{\Struct{M}}$؛ نو د لمې د فرض له مخې
$\Struct{N} \models_L !E$ هم۔ برعکس، که
$\Struct{N} \models_L !E$ وي، نو~$!D_\Struct{N}$ په~$!D$ کښې يوه
بېلېدونکې برخه ده، او دا چې~$\Struct{N} \models !D_\Struct{N}$، نو
$\Struct{N} \models_L !D$ هم۔
\end{proof}

\begin{thm}[د ليندستروم قضيه]
  \ollabel{thm:lindstrom} فرض کړئ~$\tuple{L, \models_L}$ يو نورمال
  انتزاعي منطق دے چې د متناهي شاهد او لوېنهايم--سکولم خاصيتونه لري۔
  بيا~$\tuple{L, \models_L} \le \tuple{F, \models}$؛ نو
  $\tuple{L, \models_L}$ د لومړۍ درجې له منطق سره هم‌ارز دے۔
  \textbf{د ژباړې د نسخې سمون (OLMOD-029):} سرچينې د قضيې په وينا
  کښې د نورمالوالي شرط پرېښے و، حال دا چې ثبوت د غځونې،
  آيزومورفيزم، نسبي کولو او نورو نورمالو خاصيتونو کار اخلي، او معکوسه
  څرګندونکې اړيکه هم له همدې شرط څخه راځي۔
\end{thm}

\begin{proof}
د~\olref{lem:lindstrom} له مخې، دومره ښودل کافي دي چې د هر
$!E \in L(\Lang{L})$ دپاره، چې~$\Lang{L}$ متناهي وي، داسې
$n \in \Nat$ شته چې د هر دوو !!{structure}s~$\Struct{M}$ او
$\Struct{N}$ دپاره: که~$\Struct{M} \equiv_n \Struct{N}$ وي، نو
$\Struct{M}$ او~$\Struct{N}$ پر~$!E$ سره سلا وي۔ بيا~$!E$ د لومړۍ
درجې له يوې !!{sentence} سره هم‌ارزه ده، او له دې څخه
$\tuple{L, \models_L} \le \tuple{F, \models}$ پايله کېږي۔ ځکه چې په
يوه متناهي او سوچه اړيکيزه !!{language} کښې کار کوو، د
\olref[bas][pis]{thm:b-n-f} له مخې دا وينا چې
$\Struct{M} \equiv_n \Struct{N}$ ده، په اړوندې الجبري وينا
$I_n(\emptyset,\emptyset)$ بدلولے شو۔

يو~$!E$ راکړل شي، او د تناقض دپاره فرض کړئ چې د هر~$n$ دپاره داسې
!!{structure}s~$\Struct{M}_n$ او~$\Struct{N}_n$ شته چې
$I_n(\emptyset, \emptyset)$، خو، د بېلګې په توګه،
$\Struct{M}_n \models_L !E$ او~$\Struct{N}_n \not\models_L !E$ وي۔
د آيزومورفيزم د خاصيت له مخې فرض کولے شو چې ټول~$\Struct{M}_n$ د
!!{language} ثابتونه په هماغو څيزونو تفسيروي؛ سربېره پر دې، ځکه په ژبه کښې
يوازې متناهي ډېرې اټومي !!{sentence}s شته، فرض کولے شو چې هماغه
اټومي !!{sentence}s پوره کوي؛ که داسې نۀ وي، د~$\Struct{M}$ګانو يوه
فرعي لړۍ اخيستلے شو۔ $\Struct{M}$ د ټولو~$\Struct{M}_n$ګانو اتحاد
ونيسئ؛ يعنې هغه يوازينی کوچنی !!{structure} چې هر~$\Struct{M}_n$ د
فرعي جوړښت په توګه لري۔ لکه د~\olref[lsp]{thm:abstract-p-isom} په
ثبوت کښې،~$\Struct{M}^*$ د~$\Struct{M}$ هغه غځونه ونيسئ چې !!{domain}
ئې~$\Domain{M} \cup \Domain{M}^{<\omega}$ وي، او په پراخه شوې
!!{language} کښې د نښلون محمولونه~$P$ او~$Q$ ولري۔

په ورته ډول~$\Struct{N}_n$،~$\Struct{N}$ او~$\Struct{N}^*$ تعريف
کړئ۔ اوس~$\Struct{K}$ هغه !!{structure} ونيسئ چې !!{domain} ئې د
$\Struct{M}^*$ او~$\Struct{N}^*$ د !!{domain}s تر څنګ طبيعي
عددونه~$\Nat$ د خپل طبيعي ترتيب~$\le$ سره رانغاړي۔ ژبه ئې داسې
اضافي محمولونه لري چې د اړوندو !!{domain}s~$\Domain{M}$،~$\Domain{N}$،
$\Domain{M}^{<\omega}$ او~$\Domain{N}^{<\omega}$ ښيي، او داسې
محمولونه هم لري چې د~$\Struct{M}_n$ او~$\Struct{N}_n$ د تعريف
ساحې په لاندې معنا کوډوي:
\begin{align*}
  \Domain{M_n} & = \Setabs{a \in \Domain{M}}{R(a, n)}; &
  \Domain{N_n} & = \Setabs{a \in \Domain{N}}{S(a,n)}; \\
  \Domain{M}^{<\omega}_n & = \Setabs{a \in \Domain{M}^{<\omega}}{R(a,n)}; &
  \Domain{N}^{<\omega}_n & = \Setabs{a \in \Domain{N}^{<\omega}}{S(a,n)}.
\end{align*}
!!{structure}~$\Struct{K}$ يوه درې‌ځاييزه اړيکه~$J$ هم لري، داسې چې
$J(n, \mathbf{a}, \mathbf{b})$ هله او يوازې هله صدق کوي چې
$I_n(\mathbf{a}, \mathbf{b})$ وي۔
\textbf{د ژباړې د نسخې سمون (OLMOD-030):} سرچينې دلته ګډ کوډوونکی
جوړښت د مخکېني اتحاد په هماغه نښه ښودلے، او وروسته ئې د متناهي شاهد
مدل هم د پخوانۍ غځونې په نښه ښودلے و؛ دلته ګډ جوړښت او د هغه نوی مدل
په بېلو نښو ښودل کېږي۔

اوس د~$\Lang{L}$ په هغه !!{language} کښې چې په~$R$،~$S$،~$J$ او
نورو نښو پراخه شوې ده، داسې !!a{sentence}~$!D$ شته چې وايي:~$\le$
يو منفصل خطي ترتيب دے چې لومړے خو وروستی غړے نۀ لري؛ د ترتيب
د هر~$n$ دپاره~$\Struct{M}_n \models_L !E$ او
$\Struct{N}_n \not\models_L !E$؛ او
$J(n, \mathbf{a}, \mathbf{b})$ هله او يوازې هله چې
$I_n(\mathbf{a}, \mathbf{b})$ وي۔ د انتزاعي جملې اړوندې برخې د نسبي کولو او
د کميت ټاکونکي له خاصيتونو څخه تر لاسه کېږي۔

ژبه په نويو ثابتونو~$c$،~$c_0$،~$c_1$،~\dots پراخه کړئ، او له
$!D$ سره د لومړۍ درجې هغه جملې يوځاے کړئ چې~$c_0$ د ترتيب لومړے
غړے ټاکي، هر~$c_{k+1}$ د~$c_k$ سملاسي جانشين ټاکي، او د هر معياري
$k \in \Nat$ دپاره~$c_k < c$ وايي۔ د دې تيورۍ هر متناهي فرعي سټ
په~$\Struct{K}$ کښې داسې پوره کېدے شي چې~$c$ د هغو متناهي ډېرو
يادو شوو جانشينانو څخه پورته يو طبيعي عدد وټاکي۔ د متناهي شاهد د
خاصيت له مخې ټوله تيوري يو مدل~$\Struct{K}^*$ لري۔ نو
$n^* = \Assign{c}{K^*}$ د هر معياري~$k$ دپاره د~$c_k$ له تفسير څخه
پورته دے، او د کوډ شوي منفصل ترتيب يو نامعياري غړے دے۔ په
ځانګړي ډول،~$\Struct{K}^*$ داسې فرعي !!{structure}s
$\Struct{M_{n^*}}$ او~$\Struct{N_{n^*}}$ لري چې
$\Struct{M_{n^*}} \models_L !E$ او
$\Struct{N_{n^*}} \not\models_L !E$۔ اوس د~$\Domain{M_{n^*}}$ او
$\Domain{N_{n^*}}$ د~$k$-توکيزو لړيو د جوړو يو سټ~$\mathcal{I}$
داسې تعريفولے شو:~$\tuple{\mathbf{a}, \mathbf{b}} \in \mathcal{I}$
هله او يوازې هله چې~$J(n^*-k, \mathbf{a}, \mathbf{b})$ وي؛ دلته~$k$
د~$\mathbf{a}$ او~$\mathbf{b}$ اوږدوالی دے۔ ځکه~$n^*$ نامعياري
دے، د هر معياري~$k$ دپاره~$n^* - k >0$؛ نو سټ~$\mathcal{I}$ د دې
خبرې شاهد دے چې~$\Struct{M_{n^*}} \simeq_p
\Struct{N_{n^*}}$۔ خو د~\olref[lsp]{thm:abstract-p-isom} له مخې
$\Struct{M_{n^*}}$ له~$\Struct{N_{n^*}}$ سره~$L$-هم‌ارز دے، او دا
تناقض دے۔
\textbf{د ژباړې د نسخې سمون (OLMOD-031):} سرچينې ادعا کړې وه چې د
متناهي شاهد خاصيت د يوازېنۍ جملې داسې مدل ورکوي چې نامعياري پړاو
ولري؛ هغه جمله پخپله داسې پړاو نۀ لازموي۔ دلته ژبه په يو پورني ثابت
او د معياري جانشينانو په نومونو غځول شوې، هر متناهي فرعي سټ په اصلي
کوډوونکي جوړښت کښې پوره کېږي، او متناهي شاهد ټولې تيورۍ ته له هر
معياري پړاو څخه پورته غړے ورکوي۔ د انتزاعي جملې د پوره کولو نښې هم
له اړوند منطق سره څرګندې شوې دي۔
\end{proof}

\end{document}
